\subsection{DcCA and Preliminary Results}

\begin{table}[t]
\centering
\small
\caption{Runtimes across modalities comparing DcCA with state-of-the-art baselines. N/A indicates unsupported stages. All times are in milliseconds (ms): preprocessing and indexing are measured over the dataset, while search (Step~2.a) and set selection (Step~2.b) are reported per query.}
\resizebox{\linewidth}{!}{
\begin{tabular}{l r l l l r r r r}
\specialrule{0.12em}{0pt}{0pt}
Method & Dim & Feature & Engine & Retrieval & Prep. & Index & Search & Set \\
\specialrule{0.12em}{0pt}{0pt}

\multicolumn{9}{c}{\textbf{Tabular} ($N$=10000)} \\
\midrule
Baseline & 20 & Raw & DuckDB & Linear Scan & 0.00 & 23.62 & 2.66 & N/A \\
KD-Tree  & 20 & Raw & SciPy & KD-Tree Search & 0.00 & 3.57 & 0.98 & N/A \\
DcCA     & 64 & Table2Vec & FAISS & k-NN Search & 14.03 & 0.87 & 3.91 & 3.03 \\

\specialrule{0.12em}{0pt}{0pt}
\multicolumn{9}{c}{\textbf{Temporal} ($N$=500)} \\
\midrule
Baseline & 50 & Raw & tslearn & DTW Scan & 0.00 & 0.00 & 2377.50 & N/A \\
SAX      & 10 & Symbolic & iSAX & Hamming Search & 0.93 & 0.00 & 0.49 & N/A \\
DcCA     & 128 & Time2Vec & FAISS & k-NN Search & 2.05 & 0.06 & 1.28 & 2.51 \\

\specialrule{0.12em}{0pt}{0pt}
\multicolumn{9}{c}{\textbf{Text} ($N$=500)} \\
\midrule
Baseline & 512 & Keyword & DuckDB-FTS & Keyword Match & 0.00 & 2.75 & 1.12 & N/A \\
BM25     & 512 & BM25 & Whoosh & BM25 Ranking & 0.00 & 29.25 & 14.67 & N/A \\
DcCA     & 384 & SBERT & FAISS & k-NN Search & 253.01 & 0.09 & 1.57 & 2.82 \\

\specialrule{0.12em}{0pt}{0pt}
\multicolumn{9}{c}{\textbf{Image} ($N$=100)} \\
\midrule
Baseline & 150528 & Raw & DuckDB-Blob & Linear Scan & 0.00 & 0.00 & 103.75 & N/A \\
dHash    & 64 & dHash & FastBit & Hamming Search & 0.09 & 0.00 & 0.09 & N/A \\
DcCA     & 512 & ResNet18 & FAISS & k-NN Search & 1921.81 & 0.74 & 3.69 & 3.20 \\

\specialrule{0.12em}{0pt}{0pt}
\multicolumn{9}{c}{\textbf{Graph} ($N$=2000)} \\
\midrule
Baseline & N/A & Adjacency & Neo4j & BFS Search & 0.00 & 6.30 & 0.17 & N/A \\
DcCA     & 128 & GCN & Neo4j-VEC & k-NN Search & 3.73 & 0.14 & 1.52 & 2.60 \\

\specialrule{0.12em}{0pt}{0pt}
\multicolumn{9}{c}{\textbf{MultiModal: Tabular + Text + Image} ($N$=100)} \\
\midrule
DcCA     & 512 & CLIP & FAISS & k-NN Search & 2265.25 & 0.14 & 2.32 & 2.82 \\

\specialrule{0.12em}{0pt}{0pt}
\end{tabular}
}
\label{tab:multimodal_comparison}
\end{table}



To demonstrate the feasibility of our vision, we implement \textbf{DcCA}, a prototype system that instantiates our query model, principles, and two-phase framework. \textbf{DcCA} supports tabular, temporal, graph, text, image, and multimodal data. This makes \textbf{DcCA} the only evaluated system that supports unified counterfactual querying over heterogeneous modalities, enabling constraint-aware retrieval (P1--P3) and set-level optimization (P4) in a single pipeline.

DcCA employs a shared embedding-based design for efficient cross-modal processing, using \texttt{Table2Vec}~\cite{zhang2019table2vec}, \texttt{Time2Vec}~\cite{kazemi2019time2vec}, \texttt{GCN}~\cite{kipf2016semi}, \texttt{SBERT}~\cite{reimers2019sentence}, and \texttt{ResNet18}~\cite{he2016deep}. For multimodal inputs, \textbf{DcCA} represents tabular, text, and image information in a unified embedding space via \texttt{CLIP}~\cite{radford2021learning}. Following our framework, Step~2.a performs cost-aware retrieval using \texttt{FAISS}~\cite{douze2025faiss} and \texttt{Neo4j-VEC} for graphs, while Step~2.b performs set-level selection under $\Omega$. All experiments are conducted on a Mac M2 Pro with 32 GB RAM.

To validate the efficiency of our design, we compare \textbf{DcCA} with representative system pipelines, including \texttt{DuckDB}~\cite{raasveldt2019duckdb} for relational filtering, \texttt{KD-Tree}~\cite{ram2019revisiting} for tabular nearest-neighbor search, \texttt{tslearn} (DTW)~\cite{tavenard2020tslearn} and \texttt{SAX}~\cite{yu2019novel} for temporal matching, \texttt{Neo4j}~\cite{neo4jVector} for graph traversal, \texttt{DuckDB-FTS}~\cite{raasveldt2019duckdb} and \texttt{BM25}~\cite{robertson2009probabilistic} for text retrieval, and both linear scan~\cite{ibrahim2017transforming} and hashing-based methods~\cite{wang2022dhash} for image retrieval. As shown in Table~\ref{tab:multimodal_comparison}, \textbf{DcCA} achieves substantial efficiency gains while supporting richer query semantics. For example, on temporal data ($N=500$), DTW-based retrieval requires 2377.50\,ms, whereas \textbf{DcCA} completes candidate retrieval in 1.28\,ms, with set-level optimization adding only 2.51\,ms. In the multimodal setting, \textbf{DcCA} supports joint tabular, text, and image counterfactual querying, completing candidate retrieval in 2.32\,ms and set-level selection in 2.82\,ms.

These results highlight two key advantages. First, \textbf{DcCA} integrates constraint-aware retrieval and cost-faithful ranking, avoiding the inefficiencies of post hoc filtering. Second, unlike modality-specific pipelines, \textbf{DcCA} supports set-level reasoning across both single-modality and multimodal inputs, producing diverse and interpretable counterfactual sets in real time. Overall, the prototype shows that counterfactual querying can be realized efficiently across heterogeneous data modalities, bridging current system capabilities with the requirements of data-centric counterfactual analysis.

\rev{Beyond runtime, we further study counterfactual \emph{quality}---collective-validity objectives, a comparison against representative counterfactual generators (FACE, DiCE, Wachter, GeCo), and approximation/scalability---together with the practical index-construction and maintenance costs; for space, these are reported in Appendix~\ref{app:eval}. In brief, objective-specific set selection improves every collective-validity objective over top-$k$ (which is best for none), and \textbf{DcCA} uniquely returns feasible, in-distribution counterfactuals whereas model-centric generators are infeasible and off-manifold.}